\section{The profile in arbitrary characteristic, and the characteristic-zero picture}

\noindent\textit{Addendum (27 September 2026).} Sections 1--9 treat $G=\Z/p^n$ in
characteristic $p$. This section records what the definition of the profile gives for
an arbitrary finite group $G$ over an arbitrary field $k$, and what it looks like in
characteristic $0$ (more generally, in characteristic prime to $|G|$), where the
mechanisms are entirely different: instead of Artin--Schreier layers, the profile is
governed by Sylow divisibility (Theorem~\ref{thm:div}) and by the roots-of-unity
defect $\ed_k(G)-\ed_{k(\mu_e)}(G)$ (Corollary~\ref{cor:rou}). Statements carry the
tags of Section~1; in addition, \S\ref{ss:prov} says explicitly which statements are
rewrites of results in the literature and which are new. As in the rest of the note,
literature is quoted without re-verification and the corresponding statements are
tagged [cited].

\subsection{Conventions}\label{ss:conv}
$k$ is any field, $G$ a finite group (constant group scheme over $k$), $F\supseteq k$
a field, $\tau\colon T\to\operatorname{Spec}F$ a $G$-torsor. The classifying
homomorphism $\varphi_\tau\colon\Gam_F\to G$ is defined up to conjugation, its image
$H=\im\varphi_\tau$ up to conjugacy, and $T\cong\coprod_{G/H}\operatorname{Spec}E$
with $E/F$ Galois with group $H$ (the \emph{splitting field} of $\tau$). We use, in
every characteristic: (B1) $\ed_k(\tau\otimes_FF')\le\ed_k(\tau)$; (B2)
$\ed_k(\tau)\ge\ed_{k'}(\tau)$ for $k\subseteq k'\subseteq F$; (B3) $\ed_k(\tau)=0$
iff $\tau$ descends to the algebraic closure of $k$ in $F$ (for $k$ algebraically
closed: iff $\tau$ is split; for general $k$ this is \emph{not} ``descends to $k$'':
$k=\Q$, $F=\Q(i)(t)$, $\tau=F(\sqrt{1+i})$ has $\ed_\Q=0$ but is not
$F(\sqrt a)$ with $a\in\Q$, since $N_{\Q(i)/\Q}(1+i)=2$ is not a square); (B4)
Lemma~4.1(i), purely transcendental base change does not change $\ed_k$. For a prime
$q$, $\ed_k(\tau;q)=\min\{\ed_k(\tau\otimes F'):F'/F\text{ finite separable},\
q\nmid[F':F]\}$ and $\ed_k(G;q)=\sup_\tau\ed_k(\tau;q)$ \cite[(2.5)]{RS}.
The generic torsor $\tau_\gen\colon\operatorname{Spec}L\to\operatorname{Spec}K$,
$L=k(V)$, $K=L^G$, is taken for a faithful representation $V$ (or any faithful
generically free $G$-variety with $k(V)/k$ regular); $\ed_k(\tau_\gen)=\ed_k(G)$
\cite{BF,Me}. The profile $\edd d_k(\tau)$ and the jump degrees $d_j(\tau)$ are as in
Definition~3.1.

\subsection{Formal properties, valid in every characteristic}

\begin{proposition}[shape]\label{prop:shape}
Let $\tau$ be a $G$-torsor over $F\supseteq k$.
\begin{enumerate}
\item[(a)] $d\mapsto\edd d_k(\tau)$ is non-increasing with values in
$\{0,\dots,\ed_k(\tau)\}$, $\edd1_k(\tau)=\ed_k(\tau)$, and
$d_{\ed(\tau)}=1\le d_{\ed(\tau)-1}\le\dots\le d_0\le|\im\varphi_\tau|$.
\item[(b)] If $[F':F]=e$ and $d\ge1$ then $\edd{de}_k(\tau)\le\edd d_k(\tau\otimes
F')\le\edd d_k(\tau)$; hence $d_j(\tau)\le[F':F]\,d_j(\tau\otimes F')$ for every
finite $F'/F$ and every $j$.
\item[(c)] $\edd d_k(\tau)=\min\{j:d_j(\tau)\le d\}$: the jump degrees determine the
profile.
\item[(d)] If $k\subseteq k'\subseteq F$ then $\edd d_{k'}(\tau)\le\edd d_k(\tau)$
for all $d$.
\item[(e)] $\edd d_k(\tau)=0$ iff some $F'/F$ of degree $\le d$ makes $\tau\otimes
F'$ descend to the algebraic closure of $k$ in $F'$; for $k$ algebraically closed, iff
$\tau$ is split by an extension of degree $\le d$.
\end{enumerate}
\stag{proved}
\end{proposition}

\begin{proof}
(a) Over the splitting field $E$, $[E:F]=|H|$, the torsor is split, so
$d_0\le|H|$; the rest is (B1). (b) is Proposition~3.2(b) verbatim; if
$[F'':F']=d_j(\tau\otimes F')$ realises $\ed(\tau\otimes F'')\le j$ then
$[F'':F]=[F':F]\,d_j(\tau\otimes F')$. (c) is a restatement, (d) is (B2) applied to
each $\tau\otimes F'$, (e) is (B3).
\end{proof}

\begin{proposition}[the last jump degree]\label{prop:d0}
Let $H=\im\varphi_\tau$ and $E$ the splitting field of $\tau$.
\begin{enumerate}
\item[(a)] $d_0(\tau)\le\min\{[F':F]:\tau\otimes F'\text{ split}\}=|H|$, with
equality if $k$ is algebraically closed.
\item[(b)] Let $k$ be perfect and $E/k$ regular. If $[F':F]<\infty$ and
$\tau\otimes F'$ descends to the algebraic closure $k_1$ of $k$ in $F'$, then
$[F':F]\ge|H|\cdot[k_1:k]$. In particular $d_0(\tau)=|H|$.
\item[(c)] For $\tau_\gen$ over a perfect $k$: $d_0(\tau_\gen)=|G|$, and
$\ed_k(\tau_\gen\otimes F')=0$ forces $[F':K]\ge|G|\,[k_1:k]$.
\end{enumerate}
\stag{proved}
\end{proposition}

\begin{proof}
(a) $T\otimes_FF'$ is split iff $E$ embeds $F$-linearly into $F'$, which forces
$[F':F]\ge|H|$; $F'=E$ works; for $k$ algebraically closed use (B3).
(b) Let $\tau\otimes F'\cong\tau_0\otimes_{k_1}F'$ with $\tau_0$ over $k_1$, of
splitting field $l/k_1$ Galois with group $H_0$. Comparing total spaces inside an
algebraic closure of $F'$ we may arrange $EF'=lF'$. Since $l/k$ is finite separable
and $E/k$ regular, $l\otimes_kE$ is a field and $[lE:E]=[l:k]=|H_0|[k_1:k]$; as
$lE\subseteq EF'$,
\[
|H_0|[k_1:k]\,|H|=[lE:E][E:F]\le[EF':F]=[lF':F'][F':F]\le|H_0|\,[F':F].
\]
(c) $E=L$ is regular over $k$ and $H=G$.
\end{proof}

\begin{remark}
Without regularity $d_0<|H|$ is possible ($k=\Q$, $\tau=\Q(t)(\sqrt2)$ is constant,
$d_0=1$). Proposition~\ref{prop:d0}(c) is the general form of Proposition~3.2(c). For
$k$ algebraically closed, $\edd d(\tau)\ge1$ for $d<|H|$: the last slot of the profile
is the only one understood in general.
\end{remark}

\begin{proposition}[base-field extension bound]\label{prop:bfe}
Let $k'/k$ be finite of degree $d$. Then $\edd d_k(\tau)\le\ed_{k'}(G)$. More
precisely, every field factor $F'$ of $F\otimes_kk'$ has $[F':F]\le d$ and
$\ed_k(\tau\otimes F')\le\ed_{k'}(\tau\otimes F')\le\ed_{k'}(G)$.
\stag{proved}
\end{proposition}

\begin{proof}
$F'\supseteq k'$ and $[F':F]\le d$. A descent $\tau\otimes F'\cong\tau_0\otimes_{F_0}
F'$ with $k'\subseteq F_0\subseteq F'$ and $\tr_{k'}F_0\le\ed_{k'}(G)$ is also a
descent over $k$, and $\tr_kF_0=\tr_{k'}F_0$ since $k'/k$ is algebraic.
\end{proof}

\begin{corollary}[roots-of-unity defect]\label{cor:rou}
Let $e$ be the exponent of $G$, $\operatorname{char}k\nmid e$, and
$d_e=[k(\mu_e):k]$. For every $G$-torsor $\tau$ over every $F\supseteq k$,
\[
\edd{[F(\mu_e):F]}_k(\tau)\le\edd{d_e}_k(\tau)\le\ed_{k(\mu_e)}(G).
\]
In particular: (i) if $G$ is abelian, $\edd{d_e}_k(\tau)\le\operatorname{rank}G$
\cite{BR}; (ii) if $G$ is a $p$-group and $\operatorname{char}k\ne p$,
$\edd{p-1}_k(\tau)\le\edd{[k(\mu_p):k]}_k(\tau)\le\ed_{k(\mu_p)}(G)$, the minimal
dimension of a faithful representation of $G$ over $k(\mu_p)$ \cite{KM}.
\stag{proved; the values are cited}
\end{corollary}

\begin{remark}
For $\tau_\gen$ and $F'=K(\mu_e)$ the torsor $\tau_\gen\otimes F'$ is the generic
torsor over $k(\mu_e)$, so $\ed_{k(\mu_e)}(\tau_\gen\otimes K(\mu_e))=\ed_{k(\mu_e)}(G)$
exactly, but whether $\ed_k$ of this torsor equals $\ed_{k(\mu_e)}(G)$ is not known to
us in general. [open]
\end{remark}

\subsection{Sylow divisibility of the jump degrees}

\begin{lemma}[$q$-closure]\label{lem:qcl}
Let $q$ be a prime, $S\le\Gam_F$ a Sylow pro-$q$ subgroup and $F^{(q)}=(F^{\rm sep})^S$.
Then $F^{(q)}$ is a union of finite subextensions of degree prime to $q$, every finite
separable extension of $F^{(q)}$ has $q$-power degree, and
$\ed_k(\tau;q)=\ed_k(\tau\otimes F^{(q)})$ \cite[(2.5)]{RS}. If $F'/F$ is finite
separable with $q\nmid[F':F]$, then there is an $F$-embedding $F'\hookrightarrow
F^{(q)}$, and hence $\ed_k(\tau\otimes F')\ge\ed_k(\tau;q)$.
\stag{cited; last statement proved}
\end{lemma}

\begin{proof}
$F'\otimes_FF^{(q)}$ is a product of finite separable extensions of $F^{(q)}$, each of
$q$-power degree, whose degrees sum to $[F':F]$; if none had degree $1$ the sum would
be divisible by $q$. So some factor is $F^{(q)}$, giving the embedding, and (B1)
gives the inequality. (For a \emph{given} embedding $F'\subseteq F^{\rm sep}$ the
compositum $F'F^{(q)}$ need not be $F^{(q)}$: $F=\Q$, $q=2$, $F'=\Q(\sqrt[3]2)$ has
$F'\otimes_\Q\Q^{(2)}\cong\Q^{(2)}\times\Q^{(2)}(\sqrt{-3})$.)
\end{proof}

\begin{lemma}[purely inseparable invariance]\label{lem:pi}
Let $\operatorname{char}k=p$, $k$ perfect, $G$ any finite group, $F'/F$ finite purely
inseparable. Then $\ed_k(\tau\otimes F')=\ed_k(\tau)$; consequently
$\edd d_k(\tau)$ is the minimum over \emph{separable} $F'/F$ of degree $\le d$.
\stag{proved}
\end{lemma}

\begin{proof}
As in Proposition~3.6, with $\sigma\colon F^{1/p^m}\to F$, $x\mapsto x^{p^m}$; the
$p^m$-power map extends to $F^{1/p^m}F^{\rm sep}\to F^{\rm sep}$ and commutes with all
automorphisms, so the induced isomorphism $\Gam_{F^{1/p^m}}\to\Gam_F$ is the canonical
one and the classifying homomorphism of $\sigma_*(\tau\otimes F^{1/p^m})$ is
$\varphi_\tau$. This recovers Proposition~3.6 without Witt vectors.
\end{proof}

\begin{proposition}\label{prop:genq}
For every prime $q$, $\ed_k(\tau_\gen;q)=\ed_k(G;q)$.
\stag{cited: \cite{RS,Me}; sketch below}
\end{proposition}

\begin{proof}[Sketch]
``$\le$'' is the definition. For ``$\ge$'' let $e=\ed_k(\tau_\gen;q)$, realised by a
prime-to-$q$ extension $K'/K$ and a descent to $K_0$, $\tr_kK_0\le e$; spread out to a
$G$-torsor $Z\to B$, $\dim B\le e$, and to a $G$-equivariant rational map
$W'\dashrightarrow Z$ from the normalisation $W'$ of $V$ in $L\otimes_KK'$, finite of
degree $[K':K]$ over $V$ and \'etale over a dense $G$-stable open $V_1\subseteq V$.
Given any torsor $\tau$ over an infinite $F$, the twist ${}^\tau V$ is an affine space
over $F$, so a free $F$-point in ${}^\tau V_1$ exists and gives an equivariant
$f\colon T\to V_1$; the pullback $T\times_VW'\to T$ is finite \'etale of degree
$[K':K]$, hence $\cong T\times_FA$ for a finite \'etale $F$-algebra $A$ of that degree.
Some field factor $F'$ of $A$ has $q\nmid[F':F]$ and $T_{F'}\to W'\dashrightarrow Z$ is
equivariant into a free $G$-variety, so $\ed_k(\tau\otimes F')\le\dim B\le e$ by
Lemma~2.3 (valid in any characteristic, Theorem~\ref{thm:geom}).
\end{proof}

\begin{theorem}[Sylow divisibility]\label{thm:div}
Let $\operatorname{char}k=0$, or $k$ perfect (then use Lemma~\ref{lem:pi}). Let
$\tau$ be a $G$-torsor over $F$.
\begin{enumerate}
\item[(a)] For every prime $q$ and every $j<\ed_k(\tau;q)$, $q$ divides $d_j(\tau)$:
every $F'/F$ with $\ed_k(\tau\otimes F')\le j$ has $q\mid[F':F]$.
\item[(b)] Every prime $q$ with $\ed_k(G;q)>j$ divides $d_j(\tau_\gen)$; hence
$d_j(\tau_\gen)\ge\prod_{q:\,\ed_k(G;q)>j}q$.
\item[(c)] For every $d\ge1$,
$\displaystyle\edd d_k(\tau_\gen)\ \ge\ \min_{1\le d'\le d}\ \max_{q\nmid d'}\ed_k(G;q)
\ \ge\ \max_{q>d}\ed_k(G;q)$ ($q$ prime).
\item[(d)] The prime-to-$q$ profile $\edd d_k{}'(\tau)=\min\{\ed_k(\tau\otimes F'):
[F':F]\le d,\ q\nmid[F':F]\}$ is non-increasing with terminal value $\ed_k(\tau;q)$,
reached at a finite prime-to-$q$ degree; for $\tau_\gen$ the terminal value is
$\ed_k(G;q)$.
\end{enumerate}
\stag{proved, granting Proposition~\ref{prop:genq}}
\end{theorem}

\begin{proof}
(a) is Lemma~\ref{lem:qcl}; (b) adds Proposition~\ref{prop:genq}; (c) a witness for
$\edd d$ has some degree $d'\le d$ and (a) applies to every $q\nmid d'$; (d) is
Lemma~\ref{lem:qcl} together with the finiteness argument of Proposition~3.7(a).
\end{proof}

\begin{remark}
(i) In the language of Theorem~4.2 and Remark~4.5, (a) is Reichstein--Scavia's
Theorem~10.1 applied to the correspondence of degree prime to $q$ furnished by a free
closed point; the proof above via the $q$-closure is the same mechanism.
(ii) Consistency with Sections 3--9: for $G=\Z/p^n$ in characteristic $p$,
$\ed_k(G;p)\le1$ and $\ed_k(G;q)=0$ for $q\ne p$, so Theorem~\ref{thm:div} only says
$p\mid d_0$, consistent with $d_0=p^n$; it says nothing about $d_1,\dots,d_{n-1}$,
which is the point of this note. In characteristic $0$ the constraint is much stronger.
\end{remark}

\begin{proposition}[$q$-groups]\label{prop:qgp}
Let $G$ be a $q$-group of order $q^n$, $\operatorname{char}k\ne q$ ($k$ perfect if
positive), $c=[k(\mu_q):k]$. Then:
(a) $\edd d_k(\tau_\gen)=\ed_k(G;q)$ for $c\le d\le q-1$, and
$\ed_k(G;q)\le\edd d_k(\tau_\gen)\le\ed_k(G)$ for $d<c$;
(b) $\ed_k(G;q)=\ed_{k(\mu_q)}(G)$ is the minimal dimension of a faithful
$k(\mu_q)$-representation \cite{KM};
(c) $q\mid d_j(\tau_\gen)$ for every $j<\ed_k(G;q)$, and $d_0=q^n$;
(d) if $\mu_q\subset k$: $\edd d(\tau_\gen)=\ed_k(G)$ for $d<q$ and all steps of the
profile occur at multiples of $q$, i.e.\ $\edd d(\tau_\gen)=\edd{q\lfloor d/q\rfloor}
(\tau_\gen)$ for $d\ge q$.
\stag{proved, modulo \cite{KM} and $\ed(\cdot;q)$ being invariant under prime-to-$q$
extensions of $k$}
\end{proposition}

\begin{proof}
Lower bounds: every degree $<q$ is prime to $q$ (Theorem~\ref{thm:div}). Upper bound
for $d\ge c$: Corollary~\ref{cor:rou} and $\ed_{k(\mu_q)}(G)=\ed_{k(\mu_q)}(G;q)=
\ed_k(G;q)$ \cite{KM,Me}. (c) is Theorem~\ref{thm:div}(b) and
Proposition~\ref{prop:d0}(c); (d) is (a)+(c) with $c=1$.
\end{proof}

Compare Proposition~3.7: in characteristic $p$ the prime-to-$p$ tail of a $p$-group
collapses to $\le1$; in characteristic $\ne q$ the prime-to-$q$ tail of a $q$-group is
\emph{constant}, equal to $\ed_k(G;q)$.

\subsection{Geometric meaning and induced torsors in every characteristic}

\begin{theorem}\label{thm:geom}
Theorem~4.2 and Corollary~4.4 hold for every finite group $G$, every field $k$ and
every characteristic: $\edd d_k(\tau)\le m$ iff some quasi-projective generically free
primitive $G$-variety $Z$ of dimension $\le m$ has a closed point of degree $\le d$ in
the free locus of ${}^\tau Z$, iff there are such $Z$, $F'/F$ of degree $\le d$ and a
$G$-equivariant morphism $T_{F'}\to Z^{\rm free}$. For $\tau_\gen$ this is a
$G$-correspondence $V\rightsquigarrow Z$ of degree $\le d$ into $Z^{\rm free}$.
\stag{proved}
\end{theorem}

\begin{proof}
The proof of Theorem~4.2 uses only: the existence of the twist ${}^\tau Z=(T\times_kZ)/G$
as a quasi-projective $F$-scheme compatible with base change and open immersions
\cite[\S3]{DR}; the identification of $F$-points of ${}^\tau Z$ with equivariant maps
$T\to Z_F$ (elementary for finite $G$); the fact that $Z^{\rm free}\to Z^{\rm free}/G$
is a $G$-torsor (SGA~1, V, finite group acting freely); and spreading out. None of
these depends on the characteristic or on $G$ being a cyclic $p$-group.
\end{proof}

\begin{lemma}[induced torsors]\label{lem:ind}
For finite groups $H\le G$ and an $H$-torsor $\tau_H$ over $F$,
$\ed_k(\Ind_H^G\tau_H)=\ed_k(\tau_H)$.
\stag{proved}
\end{lemma}

\begin{proof}
``$\le$'': descend and induce. ``$\ge$'': reduce to $T_H$ a field (if
$T_H=\Ind_{H'}^HT'$ with $T'$ a field, then $\Ind_H^G\tau_H=\Ind_{H'}^G\tau'$ and
$\ed(\tau_H)\le\ed(\tau')$). Then the argument of Proposition~3.9 goes through for
non-abelian $G$: a descent $T_0\to\operatorname{Spec}F_0$ has a connected component
$T_0'$ with stabiliser $H_0$; the components of $T_0'\otimes_{F_0}F$ are among the
$gT_H$ and are permuted transitively by $H_0$, so after conjugating $H$ we may assume
$T_H\subseteq T_0'\otimes_{F_0}F$ and $H\le H_0$. With $F_1=(T_0')^H$ and
$F_1\otimes_{F_0}F=\prod F_i$, the factor $T_0'\otimes_{F_1}F_{i_0}$ containing $T_H$
is an $H$-torsor over $F_{i_0}$ on whose components $H$ acts transitively; being
$H$-stable, $T_H$ is that factor, and $F_{i_0}=T_H^H=F$. So $\tau_H$ descends to
$F_1$, $\tr_kF_1=\tr_kF_0$.
\end{proof}

\begin{corollary}[subgroup profile]\label{cor:sub}
For every subgroup $H'\le G$: $\edd{[G:H']}_k(\tau_\gen)\le\ed_k(\tau_\gen\otimes
L^{H'})=\ed_k(H')$. Hence $d_j(\tau_\gen)\le\min\{[G:H']:\ed_k(H')\le j\}$.
\stag{proved}
\end{corollary}

\begin{proof}
$\tau_\gen\otimes L^{H'}=\Ind_{H'}^G(L/L^{H'})$, and $L/L^{H'}$ is the generic
$H'$-torsor of $V|_{H'}$ over $k$; apply Lemma~\ref{lem:ind}.
\end{proof}

This is the general form of Proposition~3.9 and of the root-field reductions of
Example~\ref{ex:Sn}.

\subsection{Examples in characteristic zero}

\begin{example}[$\Z/n$]\label{ex:cyc}
Let $\mu_n\subset k$, $\operatorname{char}k\nmid n$. The generic torsor
$K(\sqrt[n]a)/K$, $K=k(a)$, has $\ed=1$ and $d_0=n$ (Proposition~\ref{prop:d0}), so
its profile is $(1,\dots,1,0)$. \stag{proved}
\end{example}

\begin{theorem}[$(\Z/2)^n$]\label{thm:ea}
Let $k$ be algebraically closed of characteristic $\ne2$, $G=(\Z/2)^n$,
$F=K=k(a_1,\dots,a_n)$, $\tau_\gen=(a_1,\dots,a_n)\in H^1(F,\Z/2)^n$.
\begin{enumerate}
\item[(a)] $\edd{2^j}(\tau_\gen)\le n-j$ for $0\le j\le n$, and $d_0=2^n$.
\item[(b)] Every $F'/F$ of odd degree has $\ed(\tau_\gen\otimes F')=n$; all $d_j$,
$j<n$, are even.
\item[(c)] For every quadratic $F'/F$, $\ed(\tau_\gen\otimes F')\ge n-1$. Hence
$\edd2=\edd3=n-1$, $d_{n-1}=2$, and $d_{n-2}=4$ for $n\ge2$.
\item[(d)] $\edd d=1$ for $2^{n-1}\le d<2^n$. Complete profiles: $(1,0)$ for $n=1$;
$(2,1,1,0)$ for $n=2$; $(3,2,2,1,1,1,1,0)$ for $n=3$.
\end{enumerate}
\stag{proved; inputs cited}
\end{theorem}

\begin{proof}
(a) Over $F(\sqrt{a_1},\dots,\sqrt{a_j})$ the class is $(1,\dots,1,a_{j+1},\dots,a_n)$.
(b) $\ed_k(G;2)=\ed_k(G)=n$ \cite{BR,KM}; apply Theorem~\ref{thm:div}. (d) from (a),
(b), (c) and Proposition~\ref{prop:d0}.

(c) Write $(a_I)=\bigcup_{i\in I}(a_i)\in H^{|I|}(F,\Z/2)$. Suppose
$\ed(\tau_\gen\otimes F')\le n-2$ for $F'=F(\sqrt f)$. The class descends to $F_0$ with
$\tr_kF_0\le n-2$, so $H^{n-1}(F_0,\Z/2)=0$ \cite[II.4.2]{Se} and $(a_I)_{F'}=0$ for
all $|I|=n-1$. By the Arason exact sequence $H^{r-1}(F)\xrightarrow{(f)\cup}H^r(F)
\xrightarrow{\res}H^r(F(\sqrt f))$ \cite{Ar}, $(a_{[n]\smallsetminus i})\in(f)\cup
H^{n-2}(F)$ for all $i$. We show by induction on $m$ that for $F_m=k(a_1,\dots,a_m)$
there is no $g\in F_m^\times$ with $(a_{[m]\smallsetminus i})\in(g)\cup H^{m-2}(F_m)$
for all $i\in[m]$. For $m=1$: $(a_\emptyset)=1\ne0$ in $H^0$. For $m\ge2$ pass to
$\hat F=F_{m-1}((a_m))$, where $H^r(\hat F)=H^r(F_{m-1})\oplus H^{r-1}(F_{m-1})\cup(a_m)$
and units have $(u)=(\bar u)$ \cite[II.4.3]{Se}; the classes $(a_I)$ restrict
injectively (their residues are nonzero by induction on the residue maps). Write
$g=a_m^\varepsilon u$ and the given $x_i=y_i+z_i\cup(a_m)$. If $\varepsilon=1$, take
$i=m$: comparing the two summands of $(a_{[m-1]})=(a_mu)\cup x_m$ gives
$y_m=(\bar u)\cup z_m$ and $(a_{[m-1]})=(\bar u)\cup(\bar u)\cup z_m=0$ (as $-1\in
k^{\times2}$), a contradiction. If $\varepsilon=0$, for $i<m$ the second summands give
$(a_{[m-1]\smallsetminus i})=(\bar u)\cup z_i$ for all $i\in[m-1]$, contradicting the
induction hypothesis with $g=\bar u$. So no quadratic extension lowers $\ed$ by two;
with (a), (b) this gives $\edd2=\edd3=n-1$, $d_{n-1}=2$, $d_{n-2}=4$.
\end{proof}

\begin{remark}
Whether $\edd{2^j}(\tau_\gen)=n-j$ for $2\le j\le n-2$ is open; the first case is
$n=4$: is there a quartic extension of $k(a_1,\dots,a_4)$ splitting all six quaternion
algebras $(a_i,a_j)$? Non-Galois quartics have no kernel description as clean as
Arason's, which is what makes $\edd4$ hard. [open]
\end{remark}

\begin{example}[$\Z/p$ over $\Q$]\label{ex:Zp}
Let $p$ be an odd prime and $\tau$ any $\Z/p$-torsor over $F\supseteq\Q$. By
Corollary~\ref{cor:rou} and Example~\ref{ex:cyc}, $\edd{[F(\mu_p):F]}_\Q(\tau)\le1$,
so $d_1(\tau)\le p-1$ for every such torsor, although $\ed_\Q(\Z/p)\ge2$ for
$p\ge5$ \cite{Le07}: $\ed_k(G)=1$ in characteristic $0$ iff $G\hookrightarrow
\mathrm{PGL}_2(k)$, and $\Z/p\not\hookrightarrow\mathrm{PGL}_2(\Q)$ for $p\ge5$
since $\zeta_p+\zeta_p^{-1}\notin\Q$. The profile of the generic $\Z/p$-torsor over
$\Q$ thus drops to $1$ at a degree dividing $p-1$, prime to $p$, consistent with
Theorem~\ref{thm:div} since $\ed_\Q(\Z/p;p)=1$. Whether $d_1(\tau_\gen)=
[\Q(\zeta_p+\zeta_p^{-1}):\Q]=(p-1)/2$ is open.
\stag{proved from \cite{Le07}}
\end{example}

\begin{example}[symmetric groups]\label{ex:Sn}
Let $\operatorname{char}k=0$, $L=k(x_1,\dots,x_n)$, $F=K=k(s_1,\dots,s_n)$, so
$\tau_\gen$ is the generic degree-$n$ \'etale algebra. For $F_m=F(x_1,\dots,x_m)$,
$[F_m:F]=n(n-1)\cdots(n-m+1)$, Corollary~\ref{cor:sub} and (B4) give
\[
\ed_k(\tau_\gen\otimes F_m)=\ed_k(S_{n-m}),\qquad\text{so}\qquad
\edd{n(n-1)\cdots(n-m+1)}_k(\tau_\gen)\le\ed_k(S_{n-m}).
\]
With $\ed_k(S_n;q)=\lfloor n/q\rfloor$ \cite{MR}, Theorem~\ref{thm:div} gives
$\edd d(\tau_\gen)\ge\lfloor n/q_d\rfloor$ for the least prime $q_d>d$, so the profile
of the generic \'etale algebra decays no faster than $\sim n/(2d)$, whereas the
root-field bound needs degree $\sim n^m$ for a drop of $m$. For $S_6$ (with
$\ed_k(S_6)=3$, $\ed_k(S_5)=\ed_k(S_4)=2$, $\ed_k(S_3)=1$ \cite{BR}, and
$\ed_k(A_6)=3$ [cited from memory]): $\edd1=3$; $\edd d\ge2$ for $2\le d\le5$;
$\edd6\le2$, $\edd{120}\le1$, $d_0=720$; $2\mid d_2$, $6\mid d_1$; so
$d_2\in\{2,4,6\}$ and $d_1\in6\Z\cap[6,120]$. The discriminant field $L^{A_6}$ is
\emph{not} a witness for $d_2=2$, since over it $\ed=\ed_k(A_6)=3$.
\stag{proved from cited inputs}
\end{example}

\subsection{Provenance: rewrites versus new statements}\label{ss:prov}

\begin{center}
\small
\begin{tabular}{@{}p{0.34\textwidth}p{0.6\textwidth}@{}}
\toprule
Statement & Status relative to the literature\\
\midrule
Prop.~\ref{prop:shape} (shape) & Rewrite of Prop.~3.2 for arbitrary $G$, $k$; formal.\\
Prop.~\ref{prop:d0} (last jump) & New but elementary; (a) is folklore, (b) the
regularity refinement is new.\\
Prop.~\ref{prop:bfe}, Cor.~\ref{cor:rou} & New (elementary); the values quoted are
\cite{BR,KM}.\\
Lemma~\ref{lem:qcl} & Rewrite of \cite[(2.5)]{RS}; the embedding argument is standard.\\
Lemma~\ref{lem:pi} & Rewrite of Prop.~3.6 in a form valid for every finite $G$
(no Witt vectors).\\
Prop.~\ref{prop:genq} & Known: \cite{RS,Me} ($\ed_q$ of a versal object).\\
Thm.~\ref{thm:div} & The mechanism is \cite[Thm.~10.1]{RS}; the statement about jump
degrees and the bound (c) are new.\\
Prop.~\ref{prop:qgp} & New corollary of \cite{KM}.\\
Thm.~\ref{thm:geom} & Rewrite of Thm.~4.2 in full generality; ingredients
\cite{DR,RS}.\\
Lemma~\ref{lem:ind}, Cor.~\ref{cor:sub} & Extends Prop.~3.9 to non-abelian $G$;
probably known to experts, we found no reference.\\
Ex.~\ref{ex:cyc} & Folklore.\\
Thm.~\ref{thm:ea} & (a),(b),(d) are immediate from the above; (c), i.e.\
$\edd2=\edd3=n-1$ and $d_{n-2}=4$, is new (Arason sequence \cite{Ar} + residue
induction).\\
Ex.~\ref{ex:Zp} & New observation from \cite{Le07}.\\
Ex.~\ref{ex:Sn} & The identity $\ed(\tau_\gen\otimes F_m)=\ed(S_{n-m})$ is new but
easy; inputs \cite{BR,MR}.\\
\bottomrule
\end{tabular}
\end{center}

\subsection{Discussion and open questions}
In characteristic $0$ the profile of $\tau_\gen$ is squeezed between Sylow
divisibility (Theorem~\ref{thm:div}: a prime $q$ with $\ed_k(G;q)>j$ divides every
$d_i$, $i\le j$, and the prime-to-$q$ tail is constant), the roots-of-unity layer
(Corollary~\ref{cor:rou}: within degree $[k(\mu_e):k]$, prime to $|G|$ when $|G|$ is
a prime power, the profile drops to the representation-theoretic value
$\ed_{k(\mu_e)}(G)$), and the subgroup upper bounds (Corollary~\ref{cor:sub}). What is
unknown is whether non-Galois or non-subgroup extensions can beat the subgroup
witnesses. In characteristic $p$ for a $p$-group (Sections 3--9) there is no Sylow
divisibility beyond $p\mid d_0$ and no roots-of-unity layer; the whole profile lives in
$p$-power degrees and records Artin--Schreier--Witt layers. Theorem~\ref{thm:geom}
places both regimes in one framework: the profile is the minimal degree of a
$G$-correspondence into a free $G$-variety of bounded dimension, Theorem~\ref{thm:div}
isolates the part of that degree forced by \cite[Thm.~10.1]{RS}, and the rest, the
$q$-power part in characteristic $\ne q$ and the $p$-power part in characteristic $p$,
is where the profile carries information beyond $\ed$ and $\ed(\cdot;q)$.

\begin{question}
(1) For $(\Z/2)^n$ over an algebraically closed field of characteristic $\ne2$, is
$d_j(\tau_\gen)=2^{n-j}$ for all $j$ (known for $j\in\{0,1,n-1,n\}$)? For $(\Z/q)^n$
with $\mu_q\subset k$, the cyclic-kernel analogue of the Arason sequence should give
$d_{n-1}=q$ [expected].
(2) For $\Z/p$ over $\Q$, is $d_1(\tau_\gen)=(p-1)/2$? More generally, when is
Corollary~\ref{cor:rou} sharp, and can an extension not of the form $Fk'$ do better
than adjoining roots of unity?
(3) Determine $d_2(\tau_\gen(S_6))\in\{2,4,6\}$; does $d_j(\tau_\gen(S_n))$ grow
linearly or super-polynomially in $n$?
(4) Is the profile of $\tau_\gen(V)$ independent of the faithful representation $V$?
(Presumably yes, by specialisation as in \cite{RS}.)
\end{question}
